% گزارش Part A — آزمون اشتقاق هندسی درجهٔ ادم ماکولا
% همهٔ عددها از docs/generated/idrid_derivation_gate.json می‌آیند و با اسکریپت بازتولید می‌شوند.
% معیار پذیرش پیش از محاسبهٔ هر عدد در docs/PARTA_preregistration.md ثبت و کامیت شد.

\section{آزمون اشتقاق هندسی درجهٔ ادم ماکولا}
\label{sec:parta}

\subsection{پرسش و انگیزه}

درجهٔ ادم ماکولای دیابتی در مجموعهٔ \lr{IDRiD} تعریفی هندسی و قطعی دارد: درجهٔ صفر یعنی نبودِ
ترشح سخت؛ درجهٔ دو یعنی کوتاه‌ترین فاصلهٔ هر پیکسل ترشح سخت تا مرکز فووه‌آ کمتر یا مساوی یک قطر
دیسک نوری باشد؛ و هر حالت دیگر درجهٔ یک است.

اگر این تعریف واقعاً همان تابعی باشد که متخصص به کار می‌برد، آن‌گاه هر مجموعه‌ای که ماسک ترشح
سخت و مختصات فووه‌آ و مرز دیسک نوری داشته باشد، به یک منبع برچسب سه‌کلاسه تبدیل می‌شود. این
تنها مسیری است که در این پژوهش برای افزایش دادهٔ درجهٔ میانی یافتیم، و درجهٔ میانی همان کلاسی
است که کمبودش سقف کارایی سر ادم ماکولا را تعیین می‌کند (یافتهٔ \lr{F7}).

زیرمجموعهٔ قطعه‌بندی \lr{IDRiD} شرایط ایده‌آل این آزمون را فراهم می‌کند: ماسک پیکسلی ترشح سخت،
مختصات فووه‌آ از حاشیه‌نویسی متخصص، ماسک دیسک نوری، و مهم‌تر از همه، تفکیک ترشح سخت از ترشح نرم.
این تفکیک را پیش از آزمون بررسی کردیم و تأیید شد: ماسک‌ها در دو پوشهٔ جدا نگه‌داری می‌شوند،
۸۱ تصویر ماسک ترشح سخت دارند و ۴۰ تصویر ماسک ترشح نرم.

\subsection{معیار پذیرش، پیش از محاسبه}

معیار پیش از محاسبهٔ هر عدد نوشته و کامیت شد. تصمیم گرفتیم پنج چیز را بدون استثنا گزارش کنیم:
نرخ تطابق دقیق، ماتریس درهم‌ریختگی سه‌در‌سه، بازخوانی هر کلاس با \textbf{گزارش جداگانهٔ درجهٔ یک}،
تعداد تصویرهایی که آزمون واقعاً پوشش می‌دهد، و تکرار کل آزمون زیر سه تعریف مختلف از «یک قطر دیسک».

عمداً هیچ آستانهٔ عددی برای «قبولی» تعیین نکردیم. تعیین آستانه‌ای مانند ۸۰ درصد بحث را به
درستیِ خودِ آستانه منتقل می‌کند. به جای آن تعهد کردیم که سطر درجهٔ یک حکم را حمل کند، چون
اشتقاقی که نتواند درجهٔ میانی را جای‌گذاری کند برای هدفی که این کار را برانگیخته بی‌فایده است.

\subsection{دو محدودیت که آزمون را پیش از اجرا تضعیف کردند}

\subsubsection{مجموعهٔ \lr{IDRiD} هیچ نگاشتی میان شماره‌گذاری‌هایش منتشر نکرده است}

زیرمجموعهٔ قطعه‌بندی با \lr{IDRiD\_01} تا \lr{IDRiD\_81} شماره‌گذاری شده و زیرمجموعهٔ درجه‌بندی
با \lr{IDRiD\_001} تا \lr{IDRiD\_413} برای آموزش و \lr{IDRiD\_001} تا \lr{IDRiD\_103} برای آزمون.
این دو شماره‌گذاری به هم نگاشته نمی‌شوند و ناشر هیچ جدول تناظری منتشر نکرده است.

تناظر را از خودِ تصویرها بازسازی کردیم، با دو سیگنال مستقل که هر دو باید هم‌داستان می‌شدند:
همبستگی پیکسلی میان تصویر اصلی قطعه‌بندی و هر تصویر درجه‌بندی، با آستانهٔ بزرگ‌تر از ۰٫۹۹۹۹؛ و
فاصلهٔ مرکز ثقل ماسک دیسک نوری تا مختصات دیسک نوری منتشرشده در پروندهٔ مکان‌یابی، با آستانهٔ
کمتر از ۱۰۰ پیکسل.

هیچ‌کدام به تنهایی کافی نیست. تصویرهای ته‌چشم در مقیاس کلی به هم شبیه‌اند، پس دومین گزینهٔ
همبستگی در فاصلهٔ کمتر از ۰٫۰۱ قرار می‌گیرد؛ و دیسک نوری در بیشتر چشم‌ها در ناحیهٔ مشابهی است،
پس نزدیک‌ترین دیسک یکتا نیست. هم‌زمانی این دو قوی است: یک جفت غلط باید هم تطابق پیکسلی داشته
باشد و هم دیسک نوری هم‌مکان.

از ۸۱ تصویر، \textbf{۵۴ تصویر} تأیید شد و ۲۷ تصویر تأییدنشده ماند.

\subsubsection{زیرمجموعهٔ قطعه‌بندی برای «تصویرهای ارزشمند برای قطعه‌بندی» انتخاب شده است}

توزیع درجهٔ متخصص در ۵۴ تصویر قابل استفاده چنین است: \textbf{یک تصویر درجهٔ صفر، سه تصویر
درجهٔ یک، و ۵۰ تصویر درجهٔ دو}. علت روشن است: سازندگان مجموعه تصویرهایی را برای حاشیه‌نویسی
پیکسلی انتخاب کردند که ضایعه داشته باشند، و آن تصویرها تقریباً همه درجهٔ دو هستند.

پیامد این توزیع را باید صریح نوشت: \textbf{یک پیش‌بین ثابت که همیشه درجهٔ دو بگوید، ۹۲٫۶ درصد
تطابق می‌گیرد.} پس هر عدد تطابق کلی را باید در برابر این کف خواند، و آزمون در عمل تنها روی
\textbf{چهار تصویر} قدرت تفکیک دارد. خواننده نباید تصور کند اشتقاق روی ۵۴ تصویر آزموده شده است.

\subsection{نتیجه}

\begin{table}[h]
\centering
\caption{اشتقاق هندسی در برابر درجهٔ متخصص، زیر سه تعریف از یک قطر دیسک. تعداد نمونه‌ها در
پرانتز آمده است. سطر درجهٔ یک حکم را حمل می‌کند.}
\label{tab:parta}
\begin{tabular}{lcccc}
\hline
تعریف یک قطر دیسک & تطابق دقیق & درجهٔ ۰ (۱) & \textbf{درجهٔ ۱ (۳)} & درجهٔ ۲ (۵۰) \\
\hline
محور بزرگ           & ۹۶٫۳\٪ & ۰٫۰\٪ & \textbf{۶۶٫۷\٪} & ۱۰۰٫۰\٪ \\
قطر معادل مساحت     & ۹۶٫۳\٪ & ۰٫۰\٪ & \textbf{۶۶٫۷\٪} & ۱۰۰٫۰\٪ \\
محور کوچک           & ۹۶٫۳\٪ & ۰٫۰\٪ & \textbf{۶۶٫۷\٪} & ۱۰۰٫۰\٪ \\
\hline
پیش‌بین ثابت «درجهٔ ۲» & ۹۲٫۶\٪ & ۰٫۰\٪ & ۰٫۰\٪ & ۱۰۰٫۰\٪ \\
\hline
\end{tabular}
\end{table}

سه تعریف از قطر دیسک نتیجهٔ یکسان می‌دهند. این را پیش از آزمون به‌عنوان یک ابهام روش‌شناختی
مطرح کرده بودیم، چون دیسک نوری بیضی است و آستانه بر حسب «قطر» تعریف شده است. اکنون بسته است:
فاصلهٔ ترشح تا فووه‌آ در همهٔ نمونه‌ها از آستانه دور است — میان ۵۳ تا ۷۵۶ پیکسل در برابر قطرهای
۴۵۷ تا ۶۵۷ پیکسل — پس انتخاب تعریف حکم را عوض نمی‌کند.

\subsection{چهار تصویری که آزمون روی آن‌ها تفکیک دارد}

\begin{table}[h]
\centering
\caption{تنها تصویرهایی که درجهٔ متخصصشان دو نیست. ستون «نزدیک‌ترین لکه» اندازهٔ همان مؤلفهٔ
همبندی است که درجهٔ مشتق را تعیین می‌کند.}
\label{tab:parta_four}
\begin{tabular}{lccccc}
\hline
تصویر & ترشح سخت & نزدیک‌ترین لکه & فاصله تا فووه‌آ & مشتق & متخصص \\
\hline
\lr{IDRiD\_29} & ۲۳۴۳ پیکسل، ۵ لکه   & ۱۳۹ پیکسل & ۷۵۶ پیکسل & ۱ & ۱ \checkmark \\
\lr{IDRiD\_62} & ۴۴۴۲۷ پیکسل، ۶۸ لکه & ۶۵۷ پیکسل & ۵۶۰ پیکسل & ۱ & ۱ \checkmark \\
\lr{IDRiD\_40} & ۷۷۷۹ پیکسل، ۱۹ لکه  & \textbf{۷۴ پیکسل} & ۵۳ پیکسل & ۲ & ۱ \ding{55} \\
\lr{IDRiD\_52} & ۲۹۷۴ پیکسل، ۱۳ لکه  & \textbf{۳۴۹ پیکسل} & ۸۷ پیکسل & ۲ & ۰ \ding{55} \\
\hline
\end{tabular}
\end{table}

اشتقاق دو تصویر را درست و دو تصویر را غلط تعیین می‌کند، و هر دو خطا در جهتی است که اهمیت دارد:
تعریف می‌گوید درجهٔ دو، متخصص می‌گوید صفر یا یک.

\subsection{دو توضیح رقیب که پیش از نتیجه‌گیری رد شدند}

نتیجه‌گیری از دو تصویر ناسازگار نیازمند احتیاط است. دو توضیح رقیب وجود داشت که هر دو همین
الگو را می‌سازند، و هر دو را بررسی کردیم.

\textbf{توضیح نخست: تناظر غلط.} یک جفت اشتباه دقیقاً همین تصویر را می‌سازد — ماسکی پر از ترشح
کنار درجه‌ای بی‌ربط. برای رد آن، ماسک ترشح سخت و مختصات فووه‌آ را روی \emph{تصویر درجه‌بندی
متناظر} هم‌نهاد کردیم و پیکسل‌ها را مقایسه کردیم. میانگین قدرمطلق اختلاف برای هر چهار جفت میان
۰٫۶۵ تا ۰٫۷۵ از ۲۵۵ است، یعنی تنها نوفهٔ فشرده‌سازی \lr{JPEG}. بازرسی چشمی نیز نشان می‌دهد نشانگر
فووه‌آ روی ماکولا و دیسک نوری در سمت مورد انتظار قرار دارد. تناظر درست است.

\textbf{توضیح دوم: تعریف متفاوت.} شاید \lr{IDRiD} درجه را طور دیگری تعریف کرده باشد. متن منتشرشدهٔ
مجموعه را بررسی کردیم: تعریف دقیقاً همان است که پیاده کردیم.

\subsection{تشخیص نهایی}

پس تعریف همان است و تناظر درست است. خطا جای دیگری است، و دقیق‌تر و مفیدتر از «متخصص تعریف
دیگری به کار برده» است:

\textbf{قاعدهٔ کمینهٔ فاصله روی یک ماسک پیکسلی را کوچک‌ترین لکهٔ حاشیه‌نویسی‌شده تعیین می‌کند.}
در \lr{IDRiD\_40} لکه‌ای به اندازهٔ ۷۴ پیکسل در تصویری با ابعاد ۲۸۴۸ در ۴۲۸۸ درجه را از یک به دو
می‌برد. در هر دو خطا، لکهٔ تعیین‌کننده کوچک است و درست داخل آستانه قرار دارد؛ در هر دو مورد
درست، ترشح‌ها به‌روشنی بیرون از یک قطر دیسک‌اند.

واژه‌ای که بار معنایی را حمل می‌کند «آشکار» است. درجهٔ صفر بر پایهٔ نبودِ ترشح \emph{آشکار}
تعریف شده، یعنی آنچه بالینگر می‌بیند؛ اما ماسک قطعه‌بندی هر لکه‌ای را که یک حاشیه‌نویس در سطح
پیکسل بتواند بیابد علامت می‌زند. این‌ها دو شیء متفاوت‌اند، حاصل دو کوشش حاشیه‌نویسی جدا با دو
هدف جدا، و قاعدهٔ کمینهٔ فاصله اختلافشان را بیشینه می‌کند چون به کوچک‌ترینشان حساس است.

این یک نتیجهٔ روش‌شناختی عام است، نه ویژگی \lr{IDRiD}: \textbf{اشتقاق یک برچسب ردهای بالینی از
ماسک قطعه‌بندی با آستانهٔ فاصله، کفِ حاشیه‌نویسی ماسک را به مرز تصمیم خود تبدیل می‌کند.}

\subsection{پیامدها}

مسیر اشتقاق بسته شد. مجموعهٔ \lr{SUSTech-SYSU} ارزش مهندسی ندارد، چون دو ضعف شناخته‌شده‌اش —
ترشح‌ها به شکل جعبهٔ محصورکننده نه ماسک پیکسلی، و احتمال ادغام ترشح سخت و نرم — روی روشی سوار
می‌شود که در شرایط ایده‌آل هم شکست خورده است.

آیا آستانهٔ اندازه می‌تواند آن را ترمیم کند؟ تشخیص بالا چنین چیزی را پیشنهاد می‌کند، ولی این
داده نمی‌تواند به آن پاسخ دهد: چهار تصویر تفکیک‌کننده و دو خطا در اختیار است، پس هر آستانه‌ای
روی دو نمونه برازش می‌شود و روی هیچ نمونه‌ای اعتبارسنجی نمی‌شود. آن را به‌عنوان پرسش بعدی برای
کسی که مجموعه‌ای با توان آماری کافی دارد ثبت می‌کنیم، نه به‌عنوان چیزی که اینجا حل شود.

سرانجام، این یافته به یافتهٔ \lr{F2} سازوکار می‌دهد. پیش‌تر ثبت کرده بودیم که هیچ مجموعهٔ
سه‌کلاسهٔ دوم برای ادم ماکولا وجود ندارد. اکنون دلیلی برایش داریم: درجه از روی حاشیه‌نویسی‌هایی
که این حوزه منتشر می‌کند محاسبه‌پذیر نیست، پس ساختن چنین مجموعه‌ای به درجه‌بندی تازهٔ متخصص نیاز
دارد، نه به محاسبه روی برچسب‌های موجود.
